% GENERATED FILE. Edit canonical lessons, then run npm run generate:books.
% canonical-source-hash: fnv1a64:74ad979d362cd00b
% canonical-lessons: TA-C183-ponkal, TA-C183-uppuma, TA-C183-campar, TA-C183-catni, TA-C183-poriyal

\chapter{Pongal, Upma, Sambar, Chutney, Dry stir-fried vegetable dish}
\label{ch:ta-pongal-upma-sambar-chutney-dry-stir-fried-vegetable-dish}
\hlchaptermodality{183}

\begin{chapteropening}
\textbf{By the end of this chapter:} \emph{I can say pongal, upma, sambar, a chutney and a dry stir-fried vegetable dish.}

Complete the last lesson of chapter 183: I can say pongal, upma, sambar, a chutney and a dry stir-fried vegetable dish.
\end{chapteropening}

\section[\texorpdfstring{\ta{பொங்கல்} (poṅkal)}{பொங்கல் (poṅkal)}]{\ta{பொங்கல்} (poṅkal) — pongal, a rice and lentil dish}

\label{lesson:TA-C183-ponkal}



\begin{quote}
Before the new one: say the Tamil for a pan, then the Tamil for a lid.
\end{quote}

\subsection*{You'll want to know: \ta{பொங்கல்}}
\textbf{\ta{பொங்கல்}} — \emph{poṅkal} — \textquotedblleft{}pongal, a rice and lentil dish\textquotedblright{}.

\begin{grammarlens}[title={What you've built}]
One more word for this chapter.
\end{grammarlens}

\subsection*{Your turn}
\begin{itemize}
\raggedright
  \item \emph{Say it:} \emph{poṅkal}
  \item \emph{Say it:} \emph{poṅkal}, once more
  \item \emph{Say it:} say \emph{poṅkal}
  \item \emph{Recall:} say the Tamil for a pan, then the Tamil for a lid, then say \emph{poṅkal} again
\end{itemize}

\begin{tcolorbox}[breakable,colback=teal!4,colframe=teal!35!black,title={Before you move on}]
What does \ta{பொங்கல்} mean? (\textquotedblleft{}Pongal, a rice and lentil dish\textquotedblright{}.) Say it once more.
\end{tcolorbox}
\section[\texorpdfstring{\ta{உப்புமா} (uppumā)}{உப்புமா (uppumā)}]{\ta{உப்புமா} (uppumā) — upma, a semolina dish}

\label{lesson:TA-C183-uppuma}



\begin{quote}
Before the new one: say the Tamil for a lid, then the Tamil for pongal.
\end{quote}

\subsection*{You'll want to know: \ta{உப்புமா}}
\textbf{\ta{உப்புமா}} — \emph{uppumā} — \textquotedblleft{}upma, a semolina dish\textquotedblright{}.

\begin{grammarlens}[title={What you've built}]
One more word for this chapter.
\end{grammarlens}

\subsection*{Your turn}
\begin{itemize}
\raggedright
  \item \emph{Say it:} \emph{uppumā}
  \item \emph{Say it:} \emph{uppumā}, once more
  \item \emph{Say it:} say \emph{uppumā}
  \item \emph{Recall:} say the Tamil for a lid, then the Tamil for pongal, then say \emph{uppumā} again
\end{itemize}

\begin{tcolorbox}[breakable,colback=teal!4,colframe=teal!35!black,title={Before you move on}]
What does \ta{உப்புமா} mean? (\textquotedblleft{}Upma, a semolina dish\textquotedblright{}.) Say it once more.
\end{tcolorbox}
\section[\texorpdfstring{\ta{சாம்பார்} (cāmpār)}{சாம்பார் (cāmpār)}]{\ta{சாம்பார்} (cāmpār) — sambar, a lentil and vegetable stew}

\label{lesson:TA-C183-campar}



\begin{quote}
Before the new one: say the Tamil for pongal, then the Tamil for upma.
\end{quote}

\subsection*{You'll want to know: \ta{சாம்பார்}}
\textbf{\ta{சாம்பார்}} — \emph{cāmpār} — \textquotedblleft{}sambar, a lentil and vegetable stew\textquotedblright{}.

\begin{grammarlens}[title={What you've built}]
One more word for this chapter.
\end{grammarlens}

\subsection*{Your turn}
\begin{itemize}
\raggedright
  \item \emph{Say it:} \emph{cāmpār}
  \item \emph{Say it:} \emph{cāmpār}, once more
  \item \emph{Say it:} say \emph{cāmpār}
  \item \emph{Recall:} say the Tamil for pongal, then the Tamil for upma, then say \emph{cāmpār} again
\end{itemize}

\begin{tcolorbox}[breakable,colback=teal!4,colframe=teal!35!black,title={Before you move on}]
What does \ta{சாம்பார்} mean? (\textquotedblleft{}Sambar, a lentil and vegetable stew\textquotedblright{}.) Say it once more.
\end{tcolorbox}
\section[\texorpdfstring{\ta{சட்னி} (caṭṉi)}{சட்னி (caṭṉi)}]{\ta{சட்னி} (caṭṉi) — a chutney}

\label{lesson:TA-C183-catni}



\begin{quote}
Before the new one: say the Tamil for upma, then the Tamil for sambar.
\end{quote}

\subsection*{You'll want to know: \ta{சட்னி}}
\textbf{\ta{சட்னி}} — \emph{caṭṉi} — \textquotedblleft{}a chutney\textquotedblright{}.

\begin{grammarlens}[title={What you've built}]
One more word for this chapter.
\end{grammarlens}

\subsection*{Your turn}
\begin{itemize}
\raggedright
  \item \emph{Say it:} \emph{caṭṉi}
  \item \emph{Say it:} \emph{caṭṉi}, once more
  \item \emph{Say it:} say \emph{caṭṉi}
  \item \emph{Recall:} say the Tamil for upma, then the Tamil for sambar, then say \emph{caṭṉi} again
\end{itemize}

\begin{tcolorbox}[breakable,colback=teal!4,colframe=teal!35!black,title={Before you move on}]
What does \ta{சட்னி} mean? (\textquotedblleft{}A chutney\textquotedblright{}.) Say it once more.
\end{tcolorbox}
\section[\texorpdfstring{\ta{பொரியல்} (poriyal)}{பொரியல் (poriyal)}]{\ta{பொரியல்} (poriyal) — a dry stir-fried vegetable dish}

\label{lesson:TA-C183-poriyal}



\begin{quote}
Before the new one: say the Tamil for sambar, then the Tamil for a chutney.
\end{quote}

\subsection*{You'll want to know: \ta{பொரியல்}}
\textbf{\ta{பொரியல்}} — \emph{poriyal} — \textquotedblleft{}a dry stir-fried vegetable dish\textquotedblright{}.

\begin{grammarlens}[title={What you've built}]
One more word for this chapter.
\end{grammarlens}

\subsection*{Your turn}
\begin{itemize}
\raggedright
  \item \emph{Say it:} \emph{poriyal}
  \item \emph{Say it:} \emph{poriyal}, once more
  \item \emph{Say it:} say \emph{poriyal}
  \item \emph{Recall:} say the Tamil for sambar, then the Tamil for a chutney, then say \emph{poriyal} again
\end{itemize}

\begin{tcolorbox}[breakable,colback=teal!4,colframe=teal!35!black,title={Before you move on}]
What does \ta{பொரியல்} mean? (\textquotedblleft{}A dry stir-fried vegetable dish\textquotedblright{}.) Say it once more.
\end{tcolorbox}
